\documentclass[11pt]{article}

\usepackage[margin=1in]{geometry}
\usepackage{amsmath, amssymb}
\setlength{\parindent}{0pt}

\title{Economic Analysis 2 Notes}
\date{}
\begin{document}
\maketitle
\vspace{-1.5cm}

\section*{Profit Maximization Problem}
Intuition: Stop where $P=MC$ (must ensure maximum by SOC $<0$).
\[\max_{y}\pi(y)=py-C(y)\]
Proof (FOC): $\frac{\partial \pi}{\partial y}=p-\frac{\partial C(y)}{\partial y}=0\Rightarrow p =MC \; \text{where} \; MC=\frac{\partial C(y)}{\partial y}$.

\section*{Supply Curve}
The supply curve shows quantity produced at each price.\\
The supply curve is the MC curve itself. More precisely, the upward-sloping part of MC.

\section*{Shutdown Decision}
A firm might optimize at $y^*>0$ but still lose money. Decision exists to continue producing or consider shutting down based upon the fixed and variable costs in relation to production.
\begin{itemize}
    \item Firm should still produce where $P\geq AVC$ to minimize losses from fixed costs.
\end{itemize}
Scenario:
\begin{itemize}
    \item Produce: $y^*:\pi=Py^*-VC(y^*)-FC$
    \item Shut down: $\pi=-FC$
\end{itemize}
\section*{Producer Surplus}
Definition: Producer surplus = Revenue - Variable cost
\[ PS=Py-VC(y)\]
Graphically: Area above MC, below price.\\
Relationship to profit: $PS=\pi+FC$
\section*{Day 6 Takeaways}
\begin{itemize}
    \item Profit: $\pi=Py-C(y)$
    \item 
\end{itemize}
\section*{Competitive Equilibrium}
Definition: A competitive equilibrium is price $P^*$ and quantity $Q^*$ st. $S(P^*)=D(Q^*)$.\\

Note: In short-run equilibrium, firms can make profits and losses because EQ only requires $S=D$. Firms might be above or below break-even point. In the long run, the market adjusts (firms enter/exit, supply/demand curves shift, prices change) to balance out PnL.\\

Definition: With free entry and exit, long-run equilibrium requires: 
\begin{itemize}
    \item $P=MC$ (profit max)
    \item $P=ATC_{\text{min}}$ (zero profit)
    \item $S(D)=D(P)$ (zero sum)
\end{itemize}

Accounting Profit = Revenue - Explicit Costs\\
Economic Profit = Revenue - Opportunity Costs
\section*{Pareto Efficiency}
An allocation is Pareto efficient when you cannot make anyone better off without making someone else worse off.
\section*{Edgeworth Box}
Used to show the relationship between preferences in a two-party system.
\end{document}
